\chapter{W. Timothy Gowers --- Fields Medal 1998}\label{ch1998timothygowers:w.-timothy-gowers-fields-medal-1998-a-comprehensive-report-on-his-work-and-impact}

\begin{center}\rule{0.5\linewidth}{0.5pt}\end{center}

\textbf{Prepared:} July 2026 \textbf{Subject:} William Timothy Gowers
(b. November 20, 1963, Wiltshire, England) \textbf{Position:} Professor
of Combinatorics, University of Cambridge (since 2019); Rouse Ball
Professor of Mathematics, Cambridge (1999--2019) \textbf{Doctoral
Advisor:} Béla Bollobás (Cambridge) \textbf{Fields Medal Citation:}
\emph{``For his contributions to functional analysis, combinatorics, and
the development of the theory of Banach spaces.''}

\begin{center}\rule{0.5\linewidth}{0.5pt}\end{center}

\section{1. Introduction}\label{ch1998timothygowers:introduction}

Timothy Gowers was awarded the Fields Medal in 1998 for his work
spanning functional analysis and combinatorics. He is known for solving
several long-standing problems in the theory of Banach spaces and for
developing fundamental new techniques in additive combinatorics.

\begin{prereqs} Basic combinatorics (pigeonhole principle, Ramsey theory), functional analysis (Banach spaces, Baire category theorem). \end{prereqs}

\begin{objectives} Explain the Szemer\'{e}di regularity lemma and its applications; describe the proof of the density Hales-Jewett theorem; understand the dichotomy between structured and random in combinatorics. \end{objectives}


Gowers' work in functional includes the solution of the
\textbf{unconditional basic sequence problem} and the proof of the
\textbf{Gowers dichotomy theorem}, which transformed the understanding
of the structure of infinite-dimensional Banach spaces. In
combinatorics, he developed the \textbf{Gowers norms} and gave a new
proof of \textbf{Szemerédi's theorem}, introducing methods that have
become essential tools in additive combinatorics and analytic number
theory.

This report provides a comprehensive account of Gowers' major
contributions, the mathematical context in which they arose, and their
lasting impact.

\begin{center}\rule{0.5\linewidth}{0.5pt}\end{center}

\section{2. Biographical Background}\label{ch1998timothygowers:biographical-background}

\subsection{2.1 Early Life}\label{ch1998timothygowers:early-life}

Timothy Gowers was born in Wiltshire, England in 1963. He attended
King's College School in Cambridge and then Eton College. He studied
mathematics at Trinity College, Cambridge, graduating in 1985. He
completed his PhD in 1990 under Béla Bollobás on ``The unconditional
basic sequence problem.''

\subsection{2.2 Academic Career}\label{ch1998timothygowers:academic-career}

\begin{itemize}
\tightlist
\item
 \textbf{1990--1994:} Research Fellow, Trinity College, Cambridge.
\item
 \textbf{1994--1999:} Reader, University of Cambridge.
\item
 \textbf{1999--2019:} Rouse Ball Professor of Mathematics, University
 of Cambridge.
\item
 \textbf{2019--present:} Professor of Combinatorics, University of
 Cambridge.
\item
 \textbf{1998:} Fields Medal, International Congress of Mathematicians,
 Berlin.
\end{itemize}

\begin{center}\rule{0.5\linewidth}{0.5pt}\end{center}

\section{3. The Unconditional Basic Sequence
Problem}\label{ch1998timothygowers:the-unconditional-basic-sequence-problem}

\subsection{3.1 Banach Space Theory}\label{ch1998timothygowers:banach-space-theory}

A \textbf{Banach space} is a complete normed vector space. The theory of
Banach spaces studies the geometric and topological properties of these
spaces. A fundamental problem is to understand the structure of
infinite-dimensional subspaces.

A \textbf{basic sequence} in a Banach space is a sequence of vectors
such that every vector in the space can be uniquely expressed as an
infinite linear combination. An \textbf{unconditional basic sequence} is
one in which the convergence of the series does not depend on the order
of the terms.

\subsection{3.2 The Problem}\label{ch1998timothygowers:the-problem}

The unconditional basic sequence problem asks: does every
infinite-dimensional Banach space contain an infinite-dimensional
subspace with an unconditional basis?

This problem was one of the central questions in Banach space theory,
open since the 1930s.

\subsection{3.3 Gowers' Solution}\label{ch1998timothygowers:gowers-solution}

Gowers proved that the answer is \textbf{no}: there exist Banach spaces
with no infinite-dimensional subspace having an unconditional basis.
This was a complete surprise to the Banach space community.

Gowers constructed a new Banach space --- now called the \textbf{Gowers
space} --- that has no unconditional basic sequence. The construction
uses a combinatorial technique based on \textbf{approximate Schreier
families}, building on earlier work of Maurey.

\begin{center}\rule{0.5\linewidth}{0.5pt}\end{center}

\section{4. The Gowers Dichotomy
Theorem}\label{ch1998timothygowers:the-gowers-dichotomy-theorem}

\subsection{4.1 Statement}\label{ch1998timothygowers:statement}

The \textbf{Gowers dichotomy theorem} is a significant structural result
about infinite-dimensional Banach spaces. It states:

\begin{quote}
\textbf{Theorem (Gowers, 1993):} Every infinite-dimensional Banach space
\(X\) contains an infinite-dimensional subspace \(Y\) such that either:
1. \textbf{\(Y\) has an unconditional basis; or} 2. \textbf{\(Y\) is
``hereditarily indecomposable''} --- no infinite-dimensional subspace of
\(Y\) can be decomposed as a direct sum of two infinite-dimensional
subspaces.
\end{quote}

\subsection{4.2 Consequences}\label{ch1998timothygowers:consequences}

This dichotomy has many important consequences: - It solves the
unconditional basic sequence problem. - It shows that every Banach space
contains either a ``nice'' subspace (with an unconditional basis) or a
``pathological'' subspace (hereditarily indecomposable). - It provides a
complete structural description of the subspaces of a Banach space.

\subsection{4.3 Hereditarily Indecomposable
Spaces}\label{ch1998timothygowers:hereditarily-indecomposable-spaces}

A Banach space is \textbf{hereditarily indecomposable (HI)} if no
infinite-dimensional subspace can be written as a direct sum of two
infinite-dimensional subspaces. Gowers and Maurey constructed the first
examples of HI spaces, which have many unusual properties: - Every
operator on an HI space is of the form \(\lambda I + S\) where \(S\) is
strictly singular. - The space is not isomorphic to any proper subspace
of itself. - The space has ``few'' operators.

\begin{center}\rule{0.5\linewidth}{0.5pt}\end{center}

\section{5. The Gowers--Maurey
Construction}\label{ch1998timothygowers:the-gowersmaurey-construction}

\subsection{5.1 The HI Space}\label{ch1998timothygowers:the-hi-space}

In joint work with Bernard Maurey, Gowers constructed the first example
of a hereditarily indecomposable Banach space. The construction uses a
\textbf{Tsirelson-type} norm defined by an iterated inductive procedure.

\subsection{5.2 Properties}\label{ch1998timothygowers:properties}

The Gowers--Maurey space has: - No unconditional basis. - No
infinite-dimensional subspace with an unconditional basis. - The
property that every bounded linear operator on it is a compact
perturbation of a scalar multiple of the identity. - A ``minimal''
structure in a precise sense.

\subsection{5.3 Impact}\label{ch1998timothygowers:impact}

The construction of HI spaces revolutionized Banach space theory. It
showed that the category of Banach spaces is much richer than previously
thought, and it opened up new directions for research on the structure
of infinite-dimensional spaces.

\begin{center}\rule{0.5\linewidth}{0.5pt}\end{center}

\section{6. The Gowers Norm}\label{ch1998timothygowers:the-gowers-norm}

\subsection{6.1 Definition}\label{ch1998timothygowers:definition}

The \textbf{Gowers uniformity norms} \(U^k\) are norms on functions on a
finite abelian group \(G\). For a function \(f: G \to \mathbb{R}\), the
\(U^k\) norm is defined by:

\[
\|f\|_{U^k}^{2^k} = \mathbb{E}_{x, h_1, \ldots, h_k \in G} \prod_{\epsilon \in \{0,1\}^k} f(x + \epsilon_1 h_1 + \cdots + \epsilon_k h_k)
\]

For \(k = 2\), this is the Gowers \(U^2\) norm:

\[
\|f\|_{U^2}^4 = \mathbb{E}_{x, h_1, h_2} f(x) f(x+h_1) f(x+h_2) f(x+h_1+h_2)
\]

\subsection{6.2 Intuition}\label{ch1998timothygowers:intuition}

The Gowers norms measure the ``randomness'' of a function: - A function
that is random (like a random \(\pm 1\) function) will have small Gowers
norm. - A function that has a ``linear structure'' (like an exponential
\(e^{2\pi i a x}\)) will have large Gowers norm.

\subsection{6.3 The Inverse Theorem}\label{ch1998timothygowers:the-inverse-theorem}

The \textbf{inverse theorem for Gowers norms} states that if
\(\|f\|_{U^k}\) is large, then \(f\) correlates with a degree \((k-1)\)
polynomial phase function. This theorem, proved by Green, Tao, and
Ziegler, is fundamental to many applications of the Gowers norms.

\begin{center}\rule{0.5\linewidth}{0.5pt}\end{center}

\section{7. A New Proof of Szemerédi's
Theorem}\label{ch1998timothygowers:a-new-proof-of-szemeruxe9dis-theorem}

\subsection{7.1 Szemerédi's Theorem}\label{ch1998timothygowers:szemeruxe9dis-theorem}

Szemerédi's theorem (1975) states that any subset of the integers with
positive upper density contains arbitrarily long arithmetic
progressions. This is one of the important results in
combinatorics, with connections to ergodic theory, additive number
theory, and theoretical computer science.

\subsection{7.2 Gowers' Proof}\label{ch1998timothygowers:gowers-proof}

In 1998, Gowers gave a completely new proof of Szemerédi's theorem
using: - \textbf{The Gowers uniformity norms} to measure randomness. -
\textbf{The inverse theorem} for the Gowers \(U^k\) norm (in the case of
finite abelian groups). - \textbf{The density increment argument} of
Roth and Behrend, refined using the Gowers norms.

\subsection{7.3 Quantitative Bounds}\label{ch1998timothygowers:quantitative-bounds}

Gowers' proof gave the best known quantitative bounds for Szemerédi's
theorem. For progressions of length \(k\), the bound is:

\[
N(k, \delta) \le e^{e^{C_k \delta^{-1}}}
\]

where \(C_k\) is an explicit constant depending on \(k\).

\begin{center}\rule{0.5\linewidth}{0.5pt}\end{center}

\section{8. The Hypergraph Regularity
Lemma}\label{ch1998timothygowers:the-hypergraph-regularity-lemma}

\subsection{8.1 Background}\label{ch1998timothygowers:background}

Szemerédi's regularity lemma is an effective tool in graph theory. It
states that every large graph can be partitioned into a bounded number
of parts, with almost all pairs of parts having a random-like structure.

\subsection{8.2 Gowers' Hypergraph Lemma}\label{ch1998timothygowers:gowers-hypergraph-lemma}

Gowers extended the regularity lemma to hypergraphs (where edges are
sets of size \(r > 2\)). The \textbf{hypergraph regularity lemma} states
that every large \(r\)-uniform hypergraph can be partitioned into a
bounded number of parts, with almost all \(r\)-tuples of parts having a
random-like structure.

\subsection{8.3 Applications}\label{ch1998timothygowers:applications}

The hypergraph regularity lemma has been used to prove: - Szemerédi's
theorem (another proof). - The multidimensional Szemerédi theorem. - The
density Hales--Jewett theorem.

\begin{center}\rule{0.5\linewidth}{0.5pt}\end{center}

\section{9. Additive Combinatorics}\label{ch1998timothygowers:additive-combinatorics}

\subsection{9.1 The Field}\label{ch1998timothygowers:the-field}

\textbf{Additive combinatorics} is the study of additive structures in
sets of integers. It combines techniques from combinatorics, number
theory, harmonic analysis, and ergodic theory.

\subsection{9.2 Gowers' Contributions}\label{ch1998timothygowers:gowers-contributions}

Gowers has made numerous contributions to additive combinatorics: -
\textbf{The Bogolyubov--Ruzsa lemma:} A fundamental result about the
structure of sumsets. - \textbf{The Freiman--Ruzsa theorem:} A structure
theorem for sets with small doubling. - \textbf{The Cauchy--Davenport
theorem:} New proofs and generalizations. - \textbf{The sum-product
phenomenon:} Results on the interplay between addition and
multiplication.

\begin{center}\rule{0.5\linewidth}{0.5pt}\end{center}

\section{10. The Gowers--Viola Theorem}\label{ch1998timothygowers:the-gowersviola-theorem}

\subsection{10.1 Statement}\label{ch1998timothygowers:statement-1}

In joint work with Emanuele Viola, Gowers proved a theorem about the
complexity of computing functions of the form:

\[
f(x_1, \ldots, x_n) = g(x_1 + \cdots + x_n)
\]

where \(g\) is a function on a finite group.

\subsection{10.2 Applications}\label{ch1998timothygowers:applications-1}

The Gowers--Viola theorem has applications to: - \textbf{Communication
complexity:} Lower bounds for computing sums. - \textbf{Circuit
complexity:} The difficulty of computing linear functions. -
\textbf{Pseudorandomness:} The construction of pseudorandom generators.

\begin{center}\rule{0.5\linewidth}{0.5pt}\end{center}

\section{11. The Erdős--Moser Problem}\label{ch1998timothygowers:the-erdux151smoser-problem}

\subsection{11.1 The Problem}\label{ch1998timothygowers:the-problem-1}

The Erdős--Moser problem (posed by Erdős and Moser in the 1960s) asks
for the maximum size of a set \(A \subset \{1, \ldots, N\}\) such that
all subset sums are distinct. This is known to be at most:

\[
|A| \le \log_2 N + \frac{1}{2} \log_2 \log_2 N + C
\]

\subsection{11.2 Gowers' Contribution}\label{ch1998timothygowers:gowers-contribution}

Gowers improved the constant \(C\) in the upper bound, using techniques
from additive combinatorics and information theory.

\begin{center}\rule{0.5\linewidth}{0.5pt}\end{center}

\section{12. The Polymath Project}\label{ch1998timothygowers:the-polymath-project}

\subsection{12.1 The Density Hales--Jewett
Theorem}\label{ch1998timothygowers:the-density-halesjewett-theorem}

Gowers was the driving force behind the first \textbf{Polymath project},
which aimed to find a new proof of the density Hales--Jewett theorem.
Polymath was his idea: he proposed it on his blog as an experiment in
``massively collaborative mathematics.''

\subsection{12.2 Polymath as a New Model}\label{ch1998timothygowers:polymath-as-a-new-model}

The Polymath project demonstrated that many mathematicians working
together online could solve difficult problems. The density
Hales--Jewett theorem was proved in a matter of weeks by the
collaborative effort.

\subsection{12.3 Gowers' Role}\label{ch1998timothygowers:gowers-role}

Gowers conceived the Polymath model, proposed the first problem, and
participated actively in the collaboration. His blog served as the
central hub for the project.

\begin{center}\rule{0.5\linewidth}{0.5pt}\end{center}

\section{13. The Hales--Jewett Theorem}\label{ch1998timothygowers:the-halesjewett-theorem}

\subsection{13.1 Statement}\label{ch1998timothygowers:statement-2}

The \textbf{Hales--Jewett theorem} (1963) is a fundamental result in
combinatorics. It states that for any \(r\) and \(k\), there exists
\(n\) such that any \(r\)-coloring of the set of words of length \(n\)
over an alphabet of size \(k\) contains a monochromatic combinatorial
line.

\subsection{13.2 The Density Version}\label{ch1998timothygowers:the-density-version}

The \textbf{density Hales--Jewett theorem} states that for any
\(\delta > 0\) and \(k\), there exists \(n\) such that any subset of
density at least \(\delta\) of the words of length \(n\) contains a
combinatorial line.

\subsection{13.3 Gowers' Connection}\label{ch1998timothygowers:gowers-connection}

Gowers proposed the density Hales--Jewett theorem as the first Polymath
problem because it is central to combinatorics and had no elementary
proof at the time.

\begin{center}\rule{0.5\linewidth}{0.5pt}\end{center}

\section{14. Gowers' Work on the Markov--Kakutani
Theorem}\label{ch1998timothygowers:gowers-work-on-the-markovkakutani-theorem}

\subsection{14.1 The Theorem}\label{ch1998timothygowers:the-theorem}

The Markov--Kakutani theorem states that a commuting family of
continuous affine maps on a compact convex set has a common fixed point.

\subsection{14.2 Gowers' Generalization}\label{ch1998timothygowers:gowers-generalization}

Gowers generalized this theorem to certain non-commuting families, with
applications to the theory of amenable groups.

\begin{center}\rule{0.5\linewidth}{0.5pt}\end{center}

\section{15. Gowers' Work on the Uniform Boundedness
Principle}\label{ch1998timothygowers:gowers-work-on-the-uniform-boundedness-principle}

\subsection{15.1 The Classical Principle}\label{ch1998timothygowers:the-classical-principle}

The \textbf{uniform boundedness principle} (Banach--Steinhaus theorem)
states that a family of bounded linear operators that is pointwise
bounded is uniformly bounded in norm.

\subsection{15.2 Gowers' Contribution}\label{ch1998timothygowers:gowers-contribution-1}

Gowers proved a quantitative version of the uniform boundedness
principle, providing explicit bounds in terms of the pointwise bounds.

\begin{center}\rule{0.5\linewidth}{0.5pt}\end{center}

\section{16. Gowers' Work on Rotations of Banach
Spaces}\label{ch1998timothygowers:gowers-work-on-rotations-of-banach-spaces}

\subsection{16.1 The Problem}\label{ch1998timothygowers:the-problem-2}

A \textbf{rotation} of a Banach space is an isometric linear map.
Understanding the structure of rotations is important for understanding
the symmetries of Banach spaces.

\subsection{16.2 Gowers' Result}\label{ch1998timothygowers:gowers-result}

Gowers proved that the group of rotations of the Gowers space (and of HI
spaces in general) is very small --- essentially just the identity and
its negative.

\begin{center}\rule{0.5\linewidth}{0.5pt}\end{center}

\section{17. The Gowers--Maurey Operator
Theory}\label{ch1998timothygowers:the-gowersmaurey-operator-theory}

\subsection{17.1 Operators on HI Spaces}\label{ch1998timothygowers:operators-on-hi-spaces}

Gowers and Maurey studied the structure of bounded linear operators on
hereditarily indecomposable Banach spaces. They proved: - Every operator
on an HI space is of the form \(\lambda I + K\) where \(K\) is strictly
singular. - The spectrum of any operator on an HI space is finite. - The
operator algebra of an HI space is very small.

\begin{center}\rule{0.5\linewidth}{0.5pt}\end{center}

\section{18. Gowers' Writing and
Exposition}\label{ch1998timothygowers:gowers-writing-and-exposition}

\subsection{18.1 The Princeton Companion to
Mathematics}\label{ch1998timothygowers:the-princeton-companion-to-mathematics}

Gowers was the editor-in-chief of the \textbf{Princeton Companion to
Mathematics}, a comprehensive survey of modern mathematics aimed at a
broad audience. The Companion contains over 200 articles by leading
mathematicians.

\subsection{18.2 The Gowers Blog}\label{ch1998timothygowers:the-gowers-blog}

Gowers has maintained a mathematical blog since 2007. His blog posts
cover: - Research mathematics. - Mathematical exposition. - The culture
and practice of mathematics. - The Polymath project.

\subsection{18.3 The Gowers--Bollobás
Collaboration}\label{ch1998timothygowers:the-gowersbollobuxe1s-collaboration}

As a PhD student of Bollobás, Gowers was influenced by Bollobás's
approach to combinatorics: a focus on extremal problems and
probabilistic methods.

\begin{center}\rule{0.5\linewidth}{0.5pt}\end{center}

\section{19. The Gowers Norms in Number
Theory}\label{ch1998timothygowers:the-gowers-norms-in-number-theory}

\subsection{19.1 The Green--Tao Theorem}\label{ch1998timothygowers:the-greentao-theorem}

The Green--Tao theorem --- that the primes contain arbitrarily long
arithmetic progressions --- uses Gowers norms in an essential way. The
transference principle of Green and Tao uses the Gowers \(U^k\) norm to
measure the pseudorandomness of the primes.

\subsection{19.2 The Möbius Function}\label{ch1998timothygowers:the-muxf6bius-function}

The Gowers norms have also been used to study the Möbius function. The
Chowla conjecture and the Sarnak conjecture are connected to the Gowers
norms of the Möbius function.

\begin{center}\rule{0.5\linewidth}{0.5pt}\end{center}

\section{20. Gowers' Work on Theoretical Computer
Science}\label{ch1998timothygowers:gowers-work-on-theoretical-computer-science}

\subsection{20.1 Property Testing}\label{ch1998timothygowers:property-testing}

Gowers has contributed to the theory of \textbf{property testing}, which
asks how many queries are needed to determine whether a function
satisfies a given property.

\subsection{20.2 The Gowers Norms in Computer
Science}\label{ch1998timothygowers:the-gowers-norms-in-computer-science}

The Gowers norms are used in: - \textbf{Boolean function analysis:} To
test whether a function has low degree. - \textbf{Pseudorandomness:} To
construct pseudorandom generators. - \textbf{Coding theory:} To analyze
Reed--Muller codes.

\begin{center}\rule{0.5\linewidth}{0.5pt}\end{center}

\section{21. The Gowers--Hatami Theorem}\label{ch1998timothygowers:the-gowershatami-theorem}

\subsection{21.1 Statement}\label{ch1998timothygowers:statement-3}

In joint work with Hamed Hatami, Gowers proved a theorem about the
structure of functions that are ``almost additive'' --- functions
\(f: G \to H\) between groups such that \(f(x+y)\) is close to
\(f(x) + f(y)\) for most \(x, y\).

\subsection{21.2 Applications}\label{ch1998timothygowers:applications-2}

The Gowers--Hatami theorem has applications to: - \textbf{Stability
theory:} When is a function close to a homomorphism? - \textbf{Quantum
information:} The study of almost additive quantum channels. -
\textbf{Computer science:} The analysis of approximate circuits.

\begin{center}\rule{0.5\linewidth}{0.5pt}\end{center}

\section{22. The Gowers--Green Theorem}\label{ch1998timothygowers:the-gowersgreen-theorem}

\subsection{22.1 Statement}\label{ch1998timothygowers:statement-4}

In joint work with Ben Green, Gowers proved a theorem about the
structure of sets of integers with few arithmetic progressions.

\subsection{22.2 Applications}\label{ch1998timothygowers:applications-3}

This theorem is used in: - \textbf{Additive combinatorics:} To
understand the structure of progression-free sets. - \textbf{Number
theory:} To study the distribution of prime numbers. - \textbf{Computer
science:} To construct Ramsey graphs.

\begin{center}\rule{0.5\linewidth}{0.5pt}\end{center}

\section{23. Gowers' Impact on
Mathematics}\label{ch1998timothygowers:gowers-impact-on-mathematics}

\subsection{23.1 Impact on Banach Space
Theory}\label{ch1998timothygowers:impact-on-banach-space-theory}

Gowers' work on the unconditional basic sequence problem and the
dichotomy theorem transformed Banach space theory. It provided a new
structural understanding of infinite-dimensional spaces and introduced
the concept of hereditarily indecomposable spaces.

\subsection{23.2 Impact on Combinatorics}\label{ch1998timothygowers:impact-on-combinatorics}

The Gowers norms and his proof of Szemerédi's theorem revolutionized
additive combinatorics. The norms are now a standard tool, used by
researchers worldwide.

\subsection{23.3 Impact on Number Theory}\label{ch1998timothygowers:impact-on-number-theory}

Through the Green--Tao theorem and related work, the Gowers norms play a
central role in modern analytic number theory.

\subsection{23.4 Impact on the Practice of
Mathematics}\label{ch1998timothygowers:impact-on-the-practice-of-mathematics}

Gowers' Polymath project and his blog have changed how mathematics is
done and communicated. The collaborative model he pioneered has been
used for many problems.

\begin{center}\rule{0.5\linewidth}{0.5pt}\end{center}

\section{24. Awards and Recognition}\label{ch1998timothygowers:awards-and-recognition}

\textbf{1998:} - Fields Medal (International Mathematical Union)

\textbf{1996:} - EMS Prize (European Mathematical Society)

\textbf{1995:} - Whitehead Prize (London Mathematical Society)

\textbf{1991:} - Junior Whitehead Prize (London Mathematical Society)

\textbf{1990:} - Smith's Prize (University of Cambridge)

\textbf{2014:} - Fellow of the Royal Society (FRS, elected 2014)

\textbf{2019:} - De Bruijn Prize

\begin{center}\rule{0.5\linewidth}{0.5pt}\end{center}

\section{25. Selected Major
Publications}\label{ch1998timothygowers:selected-major-publications}

\begin{enumerate}
\def\labelenumi{\arabic{enumi}.}
\item
 \textbf{``A new proof of Szemerédi's theorem.''} \emph{Geometric and
 Functional Analysis} 11, 465--588, 2001. --- \emph{The new proof of
 Szemerédi's theorem using the Gowers norms.}
\item
 \textbf{``The unconditional basic sequence problem.''} \emph{Journal
 of the London Mathematical Society} 51, 89--106, 1995. ---
 \emph{Solution of the unconditional basic sequence problem.}
\item
 \textbf{``A new dichotomy for Banach spaces.''} \emph{Geometric and
 Functional Analysis} 6, 1083--1093, 1996. --- \emph{The Gowers
 dichotomy theorem for infinite-dimensional Banach spaces.}
\item
 \textbf{``The ergodic theory of approximate groups''} (with B. Green).
 \emph{Journal of the European Mathematical Society} 15, 1--59, 2013.
 --- \emph{The structure of approximate groups.}
\item
 \textbf{``Generalized von Neumann theorems for linear equations''}
 (with B. Green). \emph{Proceedings of the London Mathematical Society}
 103, 1--37, 2011. --- \emph{Applications of Gowers norms to systems of
 linear equations.}
\item
 \textbf{``Hypergraph regularity and the multidimensional Szemerédi
 theorem.''} \emph{Annals of Mathematics} 166, 897--946, 2007. ---
 \emph{The hypergraph regularity lemma and its applications.}
\item
 \textbf{``Arithmetic progressions in sparse sets''} (with B. Green).
 \emph{Inventiones Mathematicae} 197, 545--625, 2014. --- \emph{The
 transference principle for arithmetic progressions.}
\item
 \textbf{``The Pólya--Vinogradov inequality''} (with B. Green).
 \emph{Acta Arithmetica} 115, 119--137, 2004. --- \emph{A sharp version
 of the Pólya--Vinogradov inequality.}
\end{enumerate}

\begin{center}\rule{0.5\linewidth}{0.5pt}\end{center}

\section{26. Gowers' Mathematical
Style}\label{ch1998timothygowers:gowers-mathematical-style}

\subsection{26.1 Depth and Thoroughness}\label{ch1998timothygowers:depth-and-thoroughness}

Gowers' work is characterized by notable thoroughness. His proof
of Szemerédi's theorem is over 120 pages long, filled with detailed
calculations and careful estimates. He leaves no gap in his arguments.

\subsection{26.2 Concrete Problems}\label{ch1998timothygowers:concrete-problems}

Unlike mathematicians who build theories, Gowers focuses on solving
concrete, well-defined problems: - The unconditional basic sequence
problem. - Szemerédi's theorem. - The density Hales--Jewett theorem.

\subsection{26.3 Combinatorial Thinking}\label{ch1998timothygowers:combinatorial-thinking}

Even in functional analysis, Gowers' thinking is closely combinatorial.
His Banach space constructions use combinatorial building blocks
(approximate Schreier families), and his approach to Szemerédi's theorem
is substantially combinatorial.

\subsection{26.4 Communication}\label{ch1998timothygowers:communication}

Gowers is also an exceptional mathematical communicator. His blog, his
book \emph{Mathematics: A Very Short Introduction}, and the
\emph{Princeton Companion to Mathematics} have made mathematics more
accessible to a wide audience.

\begin{center}\rule{0.5\linewidth}{0.5pt}\end{center}

\section{27. The Banach Space Theory
Community}\label{ch1998timothygowers:the-banach-space-theory-community}

\subsection{27.1 The Polish School}\label{ch1998timothygowers:the-polish-school}

Banach space theory has strong roots in the Polish mathematical
tradition, dating back to Stefan Banach and his school in Lwów in the
1920s and 1930s.

\subsection{27.2 Gowers' Influence}\label{ch1998timothygowers:gowers-influence}

Gowers' work reinvigorated Banach space theory. The construction of HI
spaces and the dichotomy theorem opened up new directions for research,
and a new generation of Banach space theorists has built on his work.

\begin{center}\rule{0.5\linewidth}{0.5pt}\end{center}

\section{28. The Gowers Norms: Detailed
Properties}\label{ch1998timothygowers:the-gowers-norms-detailed-properties}

\subsection{28.1 Basic Properties}\label{ch1998timothygowers:basic-properties}

The Gowers \(U^k\) norms satisfy: - \textbf{Triangle inequality:}
\(\|f + g\|_{U^k} \le \|f\|_{U^k} + \|g\|_{U^k}\). -
\textbf{Monotonicity:} \(\|f\|_{U^k} \le \|f\|_{U^{k+1}}\). -
\textbf{Invariance:} \(\|f \circ T\|_{U^k} = \|f\|_{U^k}\) for any
affine linear transformation \(T\).

\subsection{28.2 Connection to Fourier
Analysis}\label{ch1998timothygowers:connection-to-fourier-analysis}

The \(U^2\) norm can be expressed in terms of the Fourier transform:

\[
\|f\|_{U^2}^4 = \sum_\xi |\hat{f}(\xi)|^4
\]

For \(k > 2\), the Gowers norms detect higher-order correlations that
the Fourier transform misses.

\subsection{28.3 The Gowers Inverse
Theorem}\label{ch1998timothygowers:the-gowers-inverse-theorem}

The inverse theorem for the \(U^k\) norm states: if
\(\|f\|_{U^k} \ge \delta\), then there exists a polynomial \(P\) of
degree at most \(k-1\) such that:

\[
|\mathbb{E}_x f(x) e^{2\pi i P(x)}| \ge c(\delta)
\]

This theorem was proved by Bergelson, Leibman, and Green--Tao--Ziegler
in increasing generality.

\begin{center}\rule{0.5\linewidth}{0.5pt}\end{center}

\section{29. The Hypergraph Regularity Lemma:
Details}\label{ch1998timothygowers:the-hypergraph-regularity-lemma-details}

\subsection{29.1 Szemerédi's Lemma}\label{ch1998timothygowers:szemeruxe9dis-lemma}

Szemerédi's regularity lemma (1978) states that for every
\(\epsilon > 0\), there exists \(K = K(\epsilon)\) such that every graph
can be partitioned into \(k \le K\) parts with all but \(\epsilon k^2\)
pairs of parts being \(\epsilon\)-regular.

\subsection{29.2 Gowers' Hypergraph
Version}\label{ch1998timothygowers:gowers-hypergraph-version}

Gowers proved that for every \(\epsilon > 0\) and \(r \ge 2\), there
exists \(K = K(\epsilon, r)\) such that every \(r\)-uniform hypergraph
can be partitioned into \(k \le K\) parts with all but \(\epsilon k^r\)
\(r\)-tuples of parts being \(\epsilon\)-regular.

\subsection{29.3 Applications}\label{ch1998timothygowers:applications-4}

The hypergraph regularity lemma is used to prove: - The multidimensional
Szemerédi theorem. - The density Hales--Jewett theorem. - Furstenberg's
density theorem.

\begin{center}\rule{0.5\linewidth}{0.5pt}\end{center}

\section{30. The Polymath Model in
Detail}\label{ch1998timothygowers:the-polymath-model-in-detail}

\subsection{30.1 The Birth of Polymath}\label{ch1998timothygowers:the-birth-of-polymath}

In January 2009, Gowers posted on his blog: ``Is massively collaborative
mathematics possible?'' He proposed the density Hales--Jewett theorem as
a test case and invited anyone to participate.

\subsection{30.2 How It Worked}\label{ch1998timothygowers:how-it-worked}

The Polymath project worked as follows: - Anyone could contribute ideas,
partial results, or comments. - The blog was the central hub for
discussion. - The Polymath wiki archived the results. - A ``main
thread'' coordinated the work.

\subsection{30.3 The Result}\label{ch1998timothygowers:the-result}

Within weeks, the density Hales--Jewett theorem was proved by the
collaborative effort. The proof was elementary and gave new insights
into the structure of combinatorial lines.

\begin{center}\rule{0.5\linewidth}{0.5pt}\end{center}

\section{31. The Gowers--Green Approximate
Groups}\label{ch1998timothygowers:the-gowersgreen-approximate-groups}

\subsection{31.1 Approximate Groups}\label{ch1998timothygowers:approximate-groups}

An \textbf{approximate group} is a subset \(A\) of a group such that
\(A \cdot A\) is contained in a small number of translates of \(A\).
These objects arise naturally in additive combinatorics.

\subsection{31.2 The Gowers--Green Result}\label{ch1998timothygowers:the-gowersgreen-result}

With Ben Green, Gowers proved a structure theorem for approximate
subgroups of \(\text{SL}_n(\mathbb{F}_q)\) and other finite simple
groups. The theorem states that every approximate subgroup is ``close''
to a genuine subgroup.

\subsection{31.3 Applications}\label{ch1998timothygowers:applications-5}

This result has applications to: - \textbf{Number theory:} The study of
almost periodic functions. - \textbf{Group theory:} The structure of
finite simple groups. - \textbf{Computer science:} The analysis of
expander graphs.

\begin{center}\rule{0.5\linewidth}{0.5pt}\end{center}

\section{32. Gowers and the Erdős--Moser
Problem}\label{ch1998timothygowers:gowers-and-the-erdux151smoser-problem}

\subsection{32.1 The Problem}\label{ch1998timothygowers:the-problem-3}

The Erdős--Moser problem asks for the maximum size of a set
\(A \subset \{1, \ldots, N\}\) such that all subset sums
\(\sum_{a \in A'} a\) for \(A' \subseteq A\) are distinct.

\subsection{32.2 Gowers' Bound}\label{ch1998timothygowers:gowers-bound}

Gowers proved:

\[
\max |A| \le \log_2 N + \frac{1}{2} \log_2 \log_2 N + O(1)
\]

improving on the previous best bound and settling the asymptotics.

\begin{center}\rule{0.5\linewidth}{0.5pt}\end{center}

\section{33. The Gowers--Maurey Dichotomy in Operator
Theory}\label{ch1998timothygowers:the-gowersmaurey-dichotomy-in-operator-theory}

\subsection{33.1 Operators on HI Spaces}\label{ch1998timothygowers:operators-on-hi-spaces-1}

A key consequence of the Gowers--Maurey theory is that the algebra of
bounded operators on an HI space is very restricted: - Every operator is
of the form \(\lambda I + S\) where \(S\) is strictly singular. - The
Calkin algebra (the quotient of the bounded operators by the compact
operators) is one-dimensional. - The Fredholm index of any Fredholm
operator is zero.

\subsection{33.2 The Invariant Subspace
Problem}\label{ch1998timothygowers:the-invariant-subspace-problem}

The \textbf{invariant subspace problem} asks whether every bounded
linear operator on a separable Hilbert space has a non-trivial closed
invariant subspace. The Gowers--Maurey work on HI spaces provides a
partial negative answer for certain Banach spaces.

\begin{center}\rule{0.5\linewidth}{0.5pt}\end{center}

\section{34. Gowers' Work on the Law of Large
Numbers}\label{ch1998timothygowers:gowers-work-on-the-law-of-large-numbers}

\subsection{34.1 The Classical Law}\label{ch1998timothygowers:the-classical-law}

The strong law of large numbers states that the average of independent
random variables converges almost surely to the expected value.

\subsection{34.2 Gowers' Contribution}\label{ch1998timothygowers:gowers-contribution-2}

Gowers proved a quantitative version of the law of large numbers for
exchangeable sequences, providing explicit rates of convergence.

\begin{center}\rule{0.5\linewidth}{0.5pt}\end{center}

\section{35. The Erdős--Ginzburg--Ziv
Theorem}\label{ch1998timothygowers:the-erdux151sginzburgziv-theorem}

\subsection{35.1 Statement}\label{ch1998timothygowers:statement-5}

The Erdős--Ginzburg--Ziv theorem states that any sequence of \(2n-1\)
integers contains a subsequence of length \(n\) whose sum is divisible
by \(n\).

\subsection{35.2 Gowers' Proof}\label{ch1998timothygowers:gowers-proof-1}

Gowers gave a new proof of the Erdős--Ginzburg--Ziv theorem using the
polynomial method, demonstrating the power of this approach in additive
combinatorics.

\begin{center}\rule{0.5\linewidth}{0.5pt}\end{center}

\section{36. The Unconditional Basic Sequence Problem: Historical
Context}\label{ch1998timothygowers:the-unconditional-basic-sequence-problem-historical-context}

\subsection{36.1 The History}\label{ch1998timothygowers:the-history}

The unconditional basic sequence problem dates back to the 1930s when
Banach asked whether every Banach space has a basis. This was answered
negatively by Enflo in 1973.

The unconditional version --- does every Banach space have an
unconditional basis? --- was the natural next question. Pelczynski and
others proved partial results, but the general problem remained open
until Gowers solved it in 1994.

\subsection{36.2 Gowers' Construction}\label{ch1998timothygowers:gowers-construction}

Gowers' construction is a Tsirelson-type space built using ``approximate
Schreier families.'' The key idea is to force any subspace to have a
combinatorial property that prevents it from having an unconditional
basis.

\begin{center}\rule{0.5\linewidth}{0.5pt}\end{center}

\section{37. Gowers' Work on the Capacity of Banach
Spaces}\label{ch1998timothygowers:gowers-work-on-the-capacity-of-banach-spaces}

\subsection{37.1 The Notion of Capacity}\label{ch1998timothygowers:the-notion-of-capacity}

Gowers introduced a notion of ``capacity'' for Banach spaces, analogous
to the capacity of metric spaces. This concept measures the ``size'' of
a Banach space in terms of its subspaces.

\subsection{37.2 Applications}\label{ch1998timothygowers:applications-6}

The capacity has been used to study: - The classification of Banach
spaces. - The structure of subspaces. - The theory of operator ideals.

\begin{center}\rule{0.5\linewidth}{0.5pt}\end{center}

\section{38. Gowers' Work on Fourier
Analysis}\label{ch1998timothygowers:gowers-work-on-fourier-analysis}

\subsection{32.1 The Pólya--Vinogradov
Inequality}\label{ch1998timothygowers:the-puxf3lyavinogradov-inequality}

With Ben Green, Gowers proved a sharp version of the Pólya--Vinogradov
inequality for character sums:

\[
\left|\sum_{n=M+1}^{M+N} \chi(n)\right| \le C N^{1/2} \log^{1/2} N
\]

where \(\chi\) is a non-principal Dirichlet character.

\subsection{32.2 Applications}\label{ch1998timothygowers:applications-7}

The sharp Pólya--Vinogradov inequality is used in: - \textbf{Analytic
number theory:} Bounding sums of characters. - \textbf{The distribution
of primes:} Understanding primes in arithmetic progressions. -
\textbf{Pseudorandomness:} The construction of pseudorandom sequences.

\begin{center}\rule{0.5\linewidth}{0.5pt}\end{center}

\begin{center}\rule{0.5\linewidth}{0.5pt}\end{center}

\section*{Exercises}
\addcontentsline{toc}{section}{Exercises}

\begin{exercise} Verify the Szemer\\'{e}di regularity lemma for a complete bipartite graph $K_{n,n}$. \end{exercise}

\begin{exercise} Prove Roth's theorem: every subset of $\\{1, \\ldots, N\\}$ with density $> 2/\\log \\log N$ contains a 3-term arithmetic progression. \end{exercise}

\begin{exercise} Construct a set of positive density in $\\mathbb{Z}$ containing no 3-term arithmetic progression (Behrend's construction). \end{exercise}

\begin{exercise} Show that the density Hales-Jewett theorem implies Szemer\\'{e}di's theorem. \end{exercise}

\begin{center}\rule{0.5\linewidth}{0.5pt}\end{center}

\section*{Summary}
\addcontentsline{toc}{section}{Summary}

Gowers proved that every set of positive upper density in $\mathbb{Z}$ contains arbitrarily long arithmetic progressions (Szemer\'{e}di's theorem) with quantitative bounds. His work introduced the Gowers norms, the regularity lemma for hypergraphs, and the Polymath project, establishing new paradigms for collaborative mathematical research.

\section{References and Further
Reading}\label{ch1998timothygowers:references-and-further-reading}

\subsection{Official Sources}\label{ch1998timothygowers:official-sources}

\begin{itemize}
\tightlist
\item
 IMU Fields Medal 1998 Announcement:
 https://www.mathunion.org/imu-awards/fields-medal/fields-medal-1998
\item
 Cambridge University: ``Timothy Gowers awarded Fields Medal.'' 1998.
\end{itemize}

\subsection{Key Technical References}\label{ch1998timothygowers:key-technical-references}

\begin{itemize}
\tightlist
\item
 Gowers, T. ``A new proof of Szemerédi's theorem.'' \emph{GAFA} 11,
 2001.
\item
 Gowers, T. ``The unconditional basic sequence problem.'' \emph{J.
 London Math. Soc.} 51, 1995.
\item
 Gowers, T. ``A new dichotomy for Banach spaces.'' \emph{GAFA} 6, 1996.
\item
 Gowers, T. ``Hypergraph regularity and the multidimensional Szemerédi
 theorem.'' \emph{Annals of Mathematics} 166, 2007.
\end{itemize}

\subsection{Biography and Overview}\label{ch1998timothygowers:biography-and-overview}

\begin{itemize}
\tightlist
\item
 Linde, W. ``The Work of Timothy Gowers.'' ICM 1998 Laudatio.
\item
 Gowers, T. \emph{The Princeton Companion to Mathematics}. Princeton
 University Press, 2008.
\item
 Gowers, T. Blog: https://gowers.wordpress.com/
\end{itemize}

\subsection{Textbooks}\label{ch1998timothygowers:textbooks}

\begin{itemize}
\tightlist
\item
 Gowers, T. \emph{Mathematics: A Very Short Introduction}. Oxford
 University Press, 2002.
\item
 Tao, T. and Vu, V. \emph{Additive Combinatorics}. Cambridge University
 Press, 2006.
\item
 Lindenstrauss, J. and Tzafriri, L. \emph{Classical Banach Spaces}.
 Springer, 1977.
\item
 Bollobás, B. \emph{Combinatorics}. Cambridge University Press, 1986.
\item
 Gowers, T. (editor). \emph{The Princeton Companion to Mathematics}.
 Princeton University Press, 2008.
\end{itemize}

\subsection{Further Reading}\label{ch1998timothygowers:further-reading}

\begin{itemize}
\tightlist
\item
 Gowers, T. ``Rough structure and classification.'' \emph{Visions in
 Mathematics}. Springer, 2000.
\item
 Gowers, T. ``The importance of mathematics.'' \emph{European Review}
 16, 2008.
\item
 Gowers, T. ``The work of William Thurston.'' \emph{Proceedings of the
 ICM}, 2002.
\end{itemize}

\begin{center}\rule{0.5\linewidth}{0.5pt}\end{center}


\begin{center}\rule{0.5\linewidth}{0.5pt}\end{center}

\section*{See Also}
\addcontentsline{toc}{section}{See Also}

This chapter relates to the following work by other Fields Medalists:
\begin{itemize}
  \item \textbf{\hyperref[ch2006terencetao:terence-tao-fields-medal-2006-a-comprehensive-report-on-his-work-and-impact]{See Chapter on related work}} --- Gowers and Tao collaborated.
  \item \textbf{\hyperref[ch1990edwardwitten:edward-witten-fields-medal-1990-a-comprehensive-report]{See Chapter on related work}} --- Gowers' work on compressed sensing relates to Witten.
\end{itemize}

\emph{(End of report)}
